\section*{Задача 3. Лифт (Лёгкая)}
\textbf{Условие.} В здании установлен инновационный лифт, начинающий движение с нулевого этажа. Этаж назначения задаётся двумя кнопками: \texttt{A} прибавляет к текущему этажу фиксированное число $A$, а \texttt{B} умножает текущий этаж на фиксированное число $B$. Требуется попасть на заданный этаж $X$ за минимальное количество нажатий. Если это невозможно, вывести \texttt{0}. Иначе вывести последовательность нажатий (по одному символу \texttt{A} или \texttt{B} на строке), минимальной длины, приводящую к этажу $X$. Запросы независимы, лифт каждый раз стартует с 0.

\textbf{Формат ввода.} В первой строке три целых числа: $N$ -- количество запросов ($1 \le N < 1000$), $A$ и $B$ -- параметры кнопок ($-1000 < A, B < 1000$, $A, B \ne 0$). В следующих $N$ строках записаны целевые этажи $X$ ($-1000 < X < 1000$, $X \ne 0$).

\textbf{Формат вывода.} Для каждого запроса вывести либо \texttt{0}, либо последовательность символов \texttt{A} и \texttt{B}, каждый на новой строке. Между ответами на разные запросы разделителей не требуется.

\begin{examplebox}
  \begin{minipage}[t]{0.48\linewidth}
    \textbf{Ввод:}\\
    \begin{verbatim}
1 1 2
9
    \end{verbatim}
  \end{minipage}%
  \hfill%
  \begin{minipage}[t]{0.48\linewidth}
    \textbf{Вывод:}\\
    \begin{verbatim}
A
B
B
B
A
    \end{verbatim}
  \end{minipage}
\end{examplebox}